\chapter{Schema and API Summary}
\label{app:schema-api}

\section{Principal persisted contracts}

The persisted contracts in \cref{tab:schema-summary} define where provenance and review state survive beyond one request.

\begin{longtable}{P{39mm}P{101mm}}
\caption{Principal persisted entities and traceability fields.}\label{tab:schema-summary}\\
\toprule
Entity & Principal fields \\
\midrule
\endfirsthead
\toprule
Entity & Principal fields \\
\midrule
\endhead
Paper & identifier; title; source authors; canonical identity links; metadata/provenance; ingestion, extraction, eligibility, OCR and review states; extraction/chunk/public-index generations; publication, rights and access states; timestamps \\
Chunk & identifier; paper foreign key; page start/end; inferred section; exact text; word/token count; source hash; extraction generation; eligibility/index status; timestamp \\
Index generation/manifest & immutable generation plus current pointer; corpus snapshot/hash; ordered chunk/source hashes; provider/model/revision; dimension/normalisation; configuration/code/index hashes; role/completeness \\
RAGAnswer & question; answer; retrieved chunk/paper/citation structures; provider/model/configuration; grounding and unsupported-claim warnings; reviewer fields; timestamps \\
Thesis Recommendation & transient request profile; ranked recommendation structures; evidence/gap/suggestion boundaries; citations; provider/model; grounding/review state; timestamp \\
Paper Artifact & paper foreign key; artifact type; generated text/JSON; citations and source sections; provider/model; grounding/review state \\
Canonical Author / Alias & canonical name; alias/provenance; identity/merge and review states; publication relation \\
Topic / PaperTopic / AuthorTopic & controlled name; paper/author links; score; source/evidence JSON; review state \\
ReviewEvent & item type/ID; previous/new state; reviewer ID/name/type/role; request ID; notes; field diff; previous/event hashes; timestamp \\
\bottomrule
\end{longtable}

\section{Principal API surfaces}

\Cref{tab:api-summary} separates public/transient functions from protected history and mutation operations.

\begin{table}[htbp]
\centering
\caption{User-facing API groups and security boundary.}
\label{tab:api-summary}
\begin{tabularx}{\textwidth}{P{36mm}Y P{30mm}}
\toprule
Group & Purpose & Access \\
\midrule
papers & approved statistics, browse/detail, extraction and artifact views & public projection \\
search & projection-filtered retrieval; technical scope exists only in protected/local code & public projection \\
ask & transient public generation; technical evaluation/persisted history separated & public projection; protected/local technical \\
recommendations & transient conversational ideas plus retained structured Finder compatibility/history & public projection; protected Finder history \\
topics/authors & approved explorer relations; technical rebuild separated & public projection; protected rebuild \\
artifacts & generated output with explicit unreviewed/approved state; generation/review operations & mixed public/protected \\
evaluation & evaluation state, corpus/index freshness, and public-safe results & public read \\
authentication/admin & local sessions and service actors; queues, feature/model/account/sync controls, preview-bound mutations, events, corrections & protected reviewer/admin \\
health/readiness & liveness and database/index completeness state & public read \\
\bottomrule
\end{tabularx}
\end{table}

FastAPI's generated OpenAPI document is available only in local-demo mode. Production disables interactive docs unless an operator creates an explicit boundary. Public request models expose no technical-scope selector. Public serializers redact local paths, private reviewer notes, draft corrections, restricted source text, and technical identity hashes where appropriate.

\section{Review-state semantics}

\texttt{needs\_review} denotes unreviewed content; \texttt{ai\_reviewed} denotes declared AI formative review; \texttt{needs\_reprocess} denotes preserved content that cannot be safely reproduced under current configuration; \texttt{reviewed}, \texttt{approved}, and \texttt{rejected} require a human actor with an appropriate role. Structural grounding remains separate from every review state.

\section{Representative executable contracts}

The following abridged listings show the implementation points behind the structured Finder, conversational Idea Generator, local-model identity, preview-bound administration, review, and privacy claims. They preserve operative field names and assertions while omitting imports, defaults not relevant to the claim, and test setup. The repository paths named in each explanation are authoritative; these excerpts are explanatory, not a second implementation.

\begingroup
\lstset{aboveskip=0.35em,belowskip=0.35em}

\begin{lstlisting}[language=Python,basicstyle=\ttfamily\footnotesize,caption={Abridged Finder request schema.}]
class ExtensionFinderRequest(BaseModel):
    interests: str = Field(min_length=1, max_length=2_000)
    skills: list[str] = Field(default_factory=list, max_length=50)
    available_time: Literal["2 weeks", "1 month", "semester"]
    project_type: Literal["software prototype", "data analysis",
                          "ML experiment", "dashboard/visualization", "other"]
    data_constraints: str = Field(max_length=1_000)
    preferred_difficulty: Literal["easy", "medium", "hard"]
    top_k: int = Field(default=5, ge=1, le=10)
    retrieval_mode: Literal["keyword", "feature_hashing",
                            "dense", "hybrid"]

    @field_validator("interests")
    @classmethod
    def interests_must_not_be_blank(cls, value: str) -> str:
        if not value.strip():
            raise ValueError("interests must not be empty")
        return value.strip()
\end{lstlisting}

\noindent\textbf{Contract.} \textbf{Path:} \codepath{backend/app/intelligence/extension\_recommender.py}. \textbf{Rationale:} bound the public profile and make supported choices explicit. \textbf{Input/output:} untrusted profile JSON becomes a normalized typed request. \textbf{Failure:} blank, oversized, or out-of-enum values produce validation errors; declared skills, time, and data access remain unverified. \textbf{Trace:} RQ4 and UI-03.

\begin{lstlisting}[language=Python,basicstyle=\ttfamily\footnotesize,caption={Abridged recommendation route boundary.}]
@router.post("/extensions")
def create_extension_recommendations(request, session):
    return recommend_extensions(session, request, persist=False)

@router.get("/extensions/history")
def extension_recommendation_history(
    session,
    _actor: Annotated[AuthenticatedActor, Depends(require_reviewer)],
):
    records = session.exec(select(ThesisRecommendation)).all()
    return [serialize_recommendation(record) for record in records]
\end{lstlisting}

\noindent\textbf{Contract.} \textbf{Path:} \codepath{backend/app/api/recommendations.py}; the prefix yields \codepath{POST /api/recommendations/extensions}. \textbf{Rationale:} separate public advice from review history. \textbf{Input/output:} a validated profile returns a transient recommendation; a reviewer-authenticated read returns serialized stored records. \textbf{Failure:} public history access is rejected and public generation cannot request persistence. \textbf{Trace:} RQ4/RQ7 and UI-03, UI-06, UI-09.

\begin{lstlisting}[language=Python,basicstyle=\ttfamily\footnotesize,caption={Abridged source-aware Idea Generator orchestration.}]
retrieval = retrieve(session, retrieval_query,
                     mode="keyword", top_k=12,
                     include_text=True, scope=retrieval_scope)
answerability = assess_answerability(retrieval_query,
                                      retrieval["results"])
matched = [row for row in retrieval["results"]
           if row["chunk_id"] in answerability["relevant_chunk_ids"]]
source_rows, alias_to_result = build_source_rows(matched)
model_output, metadata, warnings = generate_strict_ollama_output(
    build_idea_prompt(request, source_rows), settings=settings)

valid_aliases = [a for a in draft.source_aliases
                 if a in alias_to_result]
idea = {
    **draft.model_dump(exclude={"source_aliases"}),
    "basis": "paper_informed" if valid_aliases
             else "general_suggestion",
    "source_chunk_ids": [alias_to_result[a]["chunk_id"]
                         for a in valid_aliases],
}
\end{lstlisting}

\noindent\textbf{Contract.} \textbf{Path:} \codepath{backend/app/intelligence/idea\_generator.py}. \textbf{Rationale:} ensure conversational ideas cannot turn an invented label into evidence. \textbf{Input/output:} bounded conversation plus an approved retrieval scope becomes one to three typed ideas, citations, warnings, and runtime provenance. \textbf{Failure:} unmatched output remains a general suggestion; invalid aliases are replaced and warned; provider/schema failures are explicit. \textbf{Trace:} RQ4/RQ7 and UI-03/UI-09.

\begin{lstlisting}[language=Python,basicstyle=\ttfamily\footnotesize,caption={Abridged strict schema and local-model identity checks.}]
for attempt in range(2):
    raw, metadata, warnings = call_strict_ollama(
        provider, repair_prompt if attempt else prompt, schema)
    try:
        output = IdeaModelOutput.model_validate_json(raw)
        metadata["structured_output_repair_attempted"] = bool(attempt)
        return output, metadata, warnings
    except ValueError:
        continue
raise IdeaGenerationFailure("ollama_invalid_output",
                            "invalid idea format twice",
                            status_code=502)

preflight_digest = verify_ollama_model_digest(
    provider.base_url, provider.model, provider.expected_digest)
response = httpx.post(f"{provider.base_url}/api/generate",
                      json=payload, timeout=provider.timeout)
identity = verify_ollama_generation_identity(
    response.json(), configured_model=provider.model,
    expected_digest=preflight_digest)
postflight_digest = verify_ollama_model_digest(
    provider.base_url, provider.model, provider.expected_digest)
if postflight_digest != preflight_digest:
    raise ValueError("model identity changed during generation")
\end{lstlisting}

\noindent\textbf{Contract.} \textbf{Path:} \codepath{backend/app/intelligence/idea\_generator.py}. \textbf{Rationale:} bind generated structures to an observed approved local model and a typed response. \textbf{Input/output:} prompt, JSON schema, and model policy yield validated output plus digest/provenance metadata. \textbf{Failure:} unavailable or drifting identity returns a retryable unavailable state; two malformed responses return a retryable bad-gateway state. A stable tag without response-reported digest remains explicitly not fully attested. \textbf{Trace:} RQ6/RQ7 and the provider-integrity requirement.

\begin{lstlisting}[language=Python,basicstyle=\ttfamily\footnotesize,caption={Abridged actor-bound preview/execute transaction.}]
preview = bulk_candidates(session, payload.item_types,
                          approval_mode=payload.approval_mode)
operation = BulkOperation(
    operation_id=str(uuid4()),
    preview_hash=preview["preview_hash"],
    request_json=payload.model_dump(),
    requested_by=actor.actor_id,
    expires_at=now() + timedelta(minutes=10))
session.add(operation)
session.commit()

operation = session.get(BulkOperation, payload.operation_id)
fresh = bulk_candidates(session,
                        operation.request_json["item_types"])
if (operation.requested_by != actor.actor_id
        or operation.expires_at <= now()
        or payload.preview_hash != operation.preview_hash
        or fresh["preview_hash"] != operation.preview_hash):
    raise HTTPException(409, "create a fresh preview")

event = create_review_event(
    session, item_type=kind, item_id=item_id,
    action="bulk_approved", previous_status=previous,
    new_status="approved", actor=actor)
session.add(event)
session.commit()
\end{lstlisting}

\noindent\textbf{Contract.} \textbf{Path:} \codepath{backend/app/api/admin\_control.py}. \textbf{Rationale:} prevent an administrator from approving a candidate set that changed after inspection. \textbf{Input/output:} an actor-scoped preview produces a ten-minute operation and deterministic hash; execution recomputes state before applying attributed events. \textbf{Failure:} missing, expired, reused, actor-mismatched, or stale previews return conflict without mutation; catch-all review excludes missing-PDF papers. \textbf{Trace:} RQ5/RQ7 and UI-06--UI-08.

\begin{lstlisting}[language=Python,basicstyle=\ttfamily\footnotesize,caption={Abridged review-event schema.}]
class ReviewEvent(SQLModel, table=True):
    review_event_id: str = Field(primary_key=True)
    item_type: str
    item_id: str
    action: str
    previous_status: str | None
    new_status: str | None
    reviewer_id: str
    reviewer_type: str
    reviewer_role: str
    request_id: str
    diff_json: dict[str, Any]
    previous_event_hash: str | None
    event_hash: str | None
    created_at: datetime
\end{lstlisting}

\noindent\textbf{Contract.} \textbf{Path:} \codepath{backend/app/models/review.py}. \textbf{Rationale:} retain actor attribution and correction history. \textbf{Input/output:} an item transition plus actor, request, diff, and prior hash becomes a canonical hash-linked event. \textbf{Failure:} database triggers reject updates/deletes; a chain mismatch is detectable, although the hash is neither a signature nor an external timestamp. \textbf{Trace:} RQ5/RQ7 and UI-07--UI-08.

\begin{lstlisting}[language=Python,basicstyle=\ttfamily\footnotesize,caption={Core assertions from the public non-persistence test.}]
answer = client.post("/api/ask", json={"question": private_q})
finder = client.post("/api/recommendations/extensions",
                     json={"interests": private_interests})
assert answer.json()["persist"] is False
assert finder.json()["persist"] is False
assert answer.headers["cache-control"] == "no-store"
assert finder.headers["cache-control"] == "no-store"
assert client.get("/api/ask/history").status_code == 401
assert client.get("/api/recommendations/extensions/history").status_code == 401
assert session.exec(select(RAGAnswer)).all() == []
assert session.exec(select(ThesisRecommendation)).all() == []
\end{lstlisting}

\noindent\textbf{Contract.} \textbf{Path:} \codepath{backend/tests/test\_security\_privacy.py}. \textbf{Rationale:} falsify accidental retention or history disclosure. \textbf{Input/output:} private-looking public Ask/Finder requests return transient, non-cacheable responses. \textbf{Failure:} any persistence, cacheable header, or anonymous history response fails the test; external browser, proxy, operating-system, and network logs remain outside it. \textbf{Trace:} RQ7 and UI-06/UI-09.
\endgroup
